\section{Dario Amodei「We Must Pace the Frontier」}
Amodei は、能力向上が安全対策を追い越しつつあるとしてフロンティアのペースを抑えるべきだと主張した。第1段として、第三者評価者へ社員相当の恒常アクセスを与えることを Anthropic が一方的に約束したと告知した。
\dailyxsource{Dario Amodei (@DarioAmodei)}{https://x.com/DarioAmodei/status/2098773920774074715}

\section{Sam Altman が pacing に同意}
Altman は Amodei 投稿を引用し、pacing は最近数週間の OpenAI 内部議論の主要テーマだったと述べた。独立評価者への employee-like access を良い案とし、OpenAI も同様に行う意向を示した。
\dailyxsource{Sam Altman (@sama)}{https://x.com/sama/status/2098811563415150910}

\section{Elon Musk「Dario is right」}
Musk は Amodei のエッセイ告知を引用し「Dario is right」と投稿した。ただし、この短文だけでは xAI として第三者評価制度へ具体的にコミットしたとは確認できない。
\dailyxsource{Elon Musk (@elonmusk)}{https://x.com/elonmusk/status/2098789109980332057}

\section{Microsoft Copilot に Grok モデルを投入}
Microsoft 365 公式は、Grok モデルを Copilot の Word、Excel、PowerPoint に投入し、まず Frontier プログラム顧客へ限定展開すると発表した。Nadella と Musk も関連投稿で導入を歓迎した。
\dailyxsource{Microsoft 365 (@Microsoft365)}{https://x.com/Microsoft365/status/2098791167185785301}

\section{報道: OpenAI の2026年IPOを先送り}
Reuters や Forbes は、Altman が Fortune インタビューで2026年のIPOは安全上の懸念から「ill-advised moment」だと述べたと報じた。本人のXでIPO日程を直接述べた投稿は本収集では確認されていない。
\dailyxsource{Reuters (@Reuters)}{https://x.com/Reuters/status/2098875573124714898}

\section{Anthropic 脅威報告の継続議論}
Anthropic の脅威報告をもとに、兵器関連ソフトウェア開発、モデル蒸留・中継、影響工作などの二次報道が窓内でも続いた。原報告自体は観測窓外で、窓内投稿はその内容を再紹介するものが中心だった。
\dailyxsource{Iran International English (@IranIntl\_En)}{https://x.com/IranIntl_En/status/2098616773746491896}

\section{NVIDIA の Anthropic IPO 関与観測}
複数アカウントが、Anthropic の大型IPO構想と NVIDIA による最大100億ドル規模のアンカー投資検討を伝えた。投稿自身が両社未確認と注記する例もあり、Grok収集では未確認情報として扱う。
\dailyxsource{Scott Melker (@scottmelker)}{https://x.com/scottmelker/status/2098882310133186592}

\section{Siri AI の9月14日ベータ開始説}
非公式Appleニュースアカウントが、Siri AI は9月14日に iOS 27 と同時に英語ベータを開始し、他言語は10月以降と投稿した。Apple公式投稿は本収集で確認されておらず、未確認情報である。
\dailyxsource{The Apple Hub (@theapplehub)}{https://x.com/theapplehub/status/2098832015303000562}

\section{DeepSeek V4.1-Flash とスローダウン論}
コミュニティでは DeepSeek 4.1 Flash の価格・速度・性能に関する主張を前提に、米ラボのスローダウン提案は中国オープンモデル対策だと解釈する投稿が見られた。性能比較も含め投稿者側の主張である。
\dailyxsource{AI for Success (@ai\_for\_success)}{https://x.com/ai_for_success/status/2098839414629859565}

\section{Grok 4.7 は週明け以降との非公式観測}
非公式アカウントが Grok 4.7 は翌週、2.1T規模の効率モデルなどと投稿した。公式発表ではなく、日程・仕様ともGrok収集では未確認情報として扱う。
\dailyxsource{goodworse (@goodworse)}{https://x.com/goodworse/status/2098797611893317967}

\section{コーディングエージェント設定へのマルウェア警告}
投稿者は、スティーラーが .claude、.cursor、.vscode などへ悪意ある設定を配置し、エージェント経由で攻撃者コマンドを実行させ得る手法が報告されたと警告した。Grok収集では未検証。
\dailyxsource{aacle (@aacle\_)}{https://x.com/aacle_/status/2098601670225826256}

\section{Musk「Another AI attack」}
Musk が画像付きで「Another AI attack」と短く投稿した。本文には対象システム、主体、被害の具体説明がなく、何を指すかはこの投稿単体では確認できない。
\dailyxsource{Elon Musk (@elonmusk)}{https://x.com/elonmusk/status/2098789559093858580}

